%include directives.fmt

\newcommand\ApiOld{%
\begin{code}
send    : forall alpha ^ sigma. alpha -> !alpha;sigma  -> sigma
recv    : forall alpha ^ sigma. alpha -> ?alpha;sigma  -> alpha *. sigma
select ellj  : brOut (elli : Si) -> Sj
branch       : brIn (elli : Si) -> variant (elli : Si)
\end{code}
}

\newcommand\ApiNew{{%
\BorrowCodeStyle
\begin{code}
send    : forall alpha. alpha -> !alpha  -> Unit
recv    : forall alpha. alpha -> ?alpha  -> alpha
select ellj  : brOut (elli : Si) -> Sj
branch       : brIn (elli : Si) -> variant (elli : Si)
\end{code}
}}

\newcommand\DefnTree{%
\begin{code}
type Tree = Leaf | Node Int Tree Tree
\end{code}
}

\newcommand\DefnTreeChan{%
\begin{code}
type TreeChan = brOut (lab Leaf: Skip, lab Node: !Int;TreeChan;!TreeChan)
\end{code}
}

\newcommand\HtmlProtocolDefn{{%
\BorrowCodeStyle
\begin{code}
type HtmlChan   = brOut (lab Text : !String, lab Elem : !String;ChildChan)
type ChildChan  = brOut (lab Done : Skip, lab Child : HtmlChan;ChildChan)
\end{code}
}}

\newcommand\CodeSendTree{%
\begin{code}
sendTree : Tree -> forall sigma. TreeChan;sigma -> sigma
sendTree Leaf          c = select (lab Leaf) c
sendTree (Node x l r)  c =
  let c1 = select (lab Node) c   in {- |c1 : !Int;TreeChan;TreeChan;sigma| -}
  let c2 = send x c1             in {- |c2 : TreeChan;TreeChan;sigma| -}         {-"\label{ln:sendTree/c2-def}"-}
  let c3 = sendTree l c2         in {- |c3 : TreeChan;sigma| -}                  {-"\label{ln:sendTree/c2-use}"-}
  let c4 = sendTree r c3         in {- |c4 : sigma| -}
  c4
\end{code}
}

\newcommand\CodeSendTreeB{%
\begin{code}
sendTreeB : Tree -> TreeChan -> Skip
sendTreeB Leaf          c = select (lab Leaf) c
sendTreeB (Node x l r)  c =
  let c = select (lab Node) c in  {- |c : !Int;TreeChan;TreeChan| -}  {-"\label{ln:sendTreeB/select}"-}
  send x (ref c) ~                {- |c : TreeChan;TreeChan| -}       {-"\label{ln:sendTreeB/send-x}"-}
  sendTreeB l (ref c) ~           {- |c : TreeChan| -}
  sendTreeB r (ref c)             {- |c : Skip| -}
\end{code}
}

%{
%format xc = "\hphantom{\mathord{\&\mkern-1mu}}" c

\newcommand\CodeRenderExample{{%
\SaveColumns{renderExample}%
\newcommand\tyi  {|xc    : !String;ChildChan;Close|}
\newcommand\tyii {|xc    : ChildChan;Close|}
\newcommand\tyiii{|xc    : HtmlChan;ChildChan;Close|}
\newcommand\tyiv {|ref c : HtmlChan;Drop|}
\newcommand\tyv  {|xc    : Acq;ChildChan;Close|}
\newcommand\printTy[1]{\rlap{\csname ty\romannumeral ##1\endcsname}\hphantom{\tyiii}}
\begin{code}
renderUser : User -> HtmlChan;Close -> Unit
renderUser user c =
  let c = select (lab Elem) c in                      {- \printTy1 -}
  send "div" (ref c) ~                                {- \printTy2 -}
  let c = select (lab Child) c in                     {- \printTy3 -}
  fork (lambda ^ () -> renderProfile user (ref c)) ~  {- \printTy4\label{ln:renderExample/transBegin}\label{ln:renderExample/fork}-}
                                                      {- \printTy5 -}
  let msgs = loadMessages user in
  renderRest msgs (acquire c)                         {-"\label{ln:renderExample/transEnd}"-}

renderRest : Messages -> ChildChan;Close -> Unit
renderRest msgs c =
  let c = select (lab Child) c in
  renderMsgArea msgs (ref c) ~
  close (select (lab Done) c)

renderProfile : User -> HtmlChan;Drop -> Unit
renderProfile user c = ^^ ... ^^ drop c
\end{code}
}}
%}

\newcommand\CodeRenderRestElab{{%
\newcommand\tyi {|c1 : HtmlChan|}
\newcommand\tyii{|c2 : ChildChan;Close|}
\newcommand\printTy[1]{\rlap{\csname ty\romannumeral ##1\endcsname}\hphantom{\tyii}}
\begin{code}
renderRest msgs c =
  let c = select (lab Child) c in

  {- |c : HtmlChan;ChildChan;Close| -}
  let c1,c2 = lsplit c in
  {- \printTy1 -}
  {- \printTy2 -}

  renderMsgArea msgs c1 ~
  close (select (lab Done) c2)
\end{code}
}}

\newcommand\CodeRenderUserElab{{
\newcommand\tyi {|c1 : HtmlChan;Drop|}
\newcommand\tyii{|c2 : Acq;ChildChan;Close|}
\newcommand\printTy[1]{\rlap{\csname ty\romannumeral ##1\endcsname}\hphantom{\tyii}}
\begin{code}
renderUser user c = ...

  {- |c : HtmlChan;ChildChan;Close| -}
  let c1,c2 = rsplit c in
  {- \printTy1 -}
  {- \printTy2 -}

  fork (lambda ^ () -> renderProfile user c1) ~
  let msgs = loadMessages user in
  renderRest msgs (acquire c2)
\end{code}
}}

\newcommand\LocalSplitReduction{\begin{reductionexample}{localSplit}
\begin{code}
nu c b1 in
  (P1  par  let c = select (lab Child) c in
            let c1,c2 = lsplit c in
            renderMsgArea msgs c1 ~
            close (select (lab Done) c2))
\end{code}
\ExampleReduction*
\begin{code}
nu c b2 in
  (P2  par  let c1,c2 = lsplit c in
            renderMsgArea msgs c1 ~
            close (select (lab Done) c2))
\end{code}
\ExampleReduction*
\begin{code}
nu (c1 ^ c2) b2 in
  (P2  par  renderMsgArea msgs c1 ~
            close (select (lab Done) c2))
\end{code}
\ExampleReduction*
\begin{code}
nu c2 b3 in
  (P3  par  close (select (lab Done) c2))
\end{code}
\ExampleReduction
\begin{code}
nu c2 d in
  (P4 ^ [wait d] par close c2)
\end{code}
\ExampleReduction
\begin{code}
(P4 ^ [*] par ^^ *)
\end{code}
\end{reductionexample}}

\newcommand\RemoteSplitReduction{\begin{reductionexample}{remoteSplit}
\begin{code}
nu c b1 in
  (P1  par  let c1,c2 = rsplit c in
            fork (lambda ^ () -> renderProfile user c1) ~
            let msgs = ... ^^ in
            renderRest msgs (acquire c2))
\end{code}
\ExampleReduction
\begin{code}
nu (c1||c2) b1 in
  (P1  par  fork (lambda ^ () -> renderProfile user c1) ~
            let msgs = ... ^^ in
            renderRest msgs (acquire c2))
\end{code}
\ExampleReduction*
\begin{code}
nu (c1||c2) b2 in
  (P2  par  renderRest msgs (acquire c2)  par drop c1)
\end{code}
\ExampleReduction
\begin{code}
nu (||c2) b2 in
  (P2  par  renderRest msgs (acquire c2)  par *)
\end{code}
\ExampleReduction
\begin{code}
nu c2 b2 in
  (P2  par  renderRest msgs c2            par *)
\end{code}
\end{reductionexample}}

\newcommand\FnRenderUser{|renderUser|}
\newcommand\FnRenderRest{|renderRest|}
\newcommand\FnRenderProfile{|renderProfile|}

\newcommand\TyUser{|User|}
\newcommand\TyHtmlChan{|HtmlChan|}
\newcommand\TyChildChan{|ChildChan|}


%%% Local Variables:
%%% mode: latex
%%% TeX-master: "../popl2027.tex"
%%% End:
